%=============================================================================!
% 12.8  OTHER ENSEMBLES                                                       !
%=============================================================================!
\section{Other ensembles}
\label{sec:sm-other}

A sample on a bench exchanges more than energy with its surroundings: it
can swell against the air, and an electrode in a cell exchanges atoms and
electrons with the electrolyte and the circuit. This section extends the
argument of \cref{sec:sm-canonical} to a bath that also takes volume or
particles, finds the free energy each such ensemble makes least, and ends
with the electrons of an electronic-structure calculation.

\subsection*{Volume and particles}

The entropy of the bath depends on its volume and number of atoms as well
as its energy, and its slopes along them, each taken with the other two
held fixed (\cref{sec:ha-equations}), define the pressure $P$ and the
chemical potential $\mu$:
\begin{equation*}
   \frac{\partial S}{\partial E} = \frac1T , \qquad
   \frac{\partial S}{\partial V} = \frac PT , \qquad
   \frac{\partial S}{\partial N} = -\frac\mu T .
\end{equation*}
For the ideal gas of \cref{sec:sm-micro}, $\Omega \propto V^N$ gives
$\partial S/\partial V = N\kB/V$, so $PV = N\kB T$, the ideal gas law;
Chapter~14 shows that this $P$ is the force per unit area on the walls.
The minus sign makes $\mu$ the energy that one more atom brings, as $\dd
E$ below shows. As \cref{sec:sm-micro} showed for energy, two systems that
exchange volume settle where their values of $P/T$ agree, and two that
exchange atoms where their values of $\mu/T$ agree, once the atoms are
counted correctly. Swapping two atoms of one kind gives back the same
arrangement, so the $N!$ orderings of $N$ identical atoms are one
microstate, and $\Omega$ and $Z$ for $N$ such atoms are divided by $N!$, a
factor that changed nothing while $N$ was fixed. Without it the ideal gas,
with $\Omega \propto V^N$, gains $\ln V$ in $\ln\Omega$ for each atom
added whatever $N$ is, and two boxes sharing atoms would not settle at
equal densities, the larger box taking all the atoms;
with it, one more atom changes $\ln(V^N/N!)$ by $\ln V - \ln(N + 1)
\approx \ln(V/N)$, and two boxes at one temperature settle at equal
densities $N/V$.

A system that takes the volume $V_s$ from the bath as well as the energy
$E_s$ leaves the bath $\Omega_{\mathrm B}(E - E_s, V - V_s)$ microstates,
and the Taylor series of \cref{sec:sm-canonical}, now in two variables
with one term for each as in \cref{sec:en-gradient}, gives
\begin{equation*}
   \ln\Omega_{\mathrm B}(E - E_s, V - V_s) \approx \ln\Omega_{\mathrm B}(E, V)
   - \frac{E_s}{\kB T} - \frac{PV_s}{\kB T} ,
\end{equation*}
so a state of the system with energy $E_s$ and volume $V_s$ has the
probability proportional to $\mathrm{e}^{-\beta(E_s + PV_s)}$: the
isothermal-isobaric, or $NPT$, ensemble. A system that takes $N_s$ atoms
from the bath has, by the same steps with $\partial\ln\Omega_{\mathrm
B}/\partial N = -\mu/\kB T$,
\begin{equation*}
   \ln\Omega_{\mathrm B}(E - E_s, N - N_s) \approx \ln\Omega_{\mathrm B}(E, N)
   - \frac{E_s}{\kB T} + \frac{\mu N_s}{\kB T} ,
\end{equation*}
and so the weight $\mathrm{e}^{-\beta(E_s - \mu N_s)}$: the grand
canonical, or $\mu VT$, ensemble. Grouping the states by their volume or
their number of atoms, the sum over each group being a canonical partition
function, gives
\begin{equation*}
   Z_{NPT} = \int\mathrm{e}^{-\beta PV}Z(N, V, T)\,\dd V , \qquad
   \Xi = \sum_N\mathrm{e}^{\beta\mu N}Z(N, V, T) ,
\end{equation*}
and the derivatives of their logarithms give averages and fluctuations as
before: the volume of an $NPT$ run fluctuates with the variance $-\kB
T\,\dd\ens{V}/\dd P$ (\cref{ex:sm-volume}), large for a soft material and
small for a stiff one, which Chapter~14 uses to test a barostat's
coupling.

\subsection*{The free energies}

For a system in contact with a bath at $T$, the total energy is divided as
in \cref{sec:sm-micro}: when the system holds $E_s$ it has $\Omega(E_s)$
microstates and the entropy $S = \kB\ln\Omega(E_s)$, the bath has
$\Omega_{\mathrm B}(E - E_s)$, and the most probable division has the most
microstates of the two together, $\Omega(E_s)\,\Omega_{\mathrm B}(E -
E_s)$. Their logarithm times $\kB$ is the total entropy, which by the
expansion of \cref{sec:sm-canonical} is
\begin{equation*}
   S + S_{\mathrm B}(E - E_s) \approx S + S_{\mathrm B}(E)
   - \frac{E_s}{T} = S_{\mathrm B}(E) - \frac{E_s - TS}{T} .
\end{equation*}
$S_{\mathrm B}(E)$ is the same for every division, so the system settles
at the energy that makes $E_s - TS$ least, and for a large system that
least value is the free energy $F$ of \cref{sec:sm-canonical}. With a bath
that also takes volume, it is $F + PV$ that is least, the Gibbs free
energy, $-\kB T\ln Z_{NPT}$; with one that takes atoms, $F - \mu N$, the
grand potential, $-\kB T\ln\Xi$. From here on $E$, $V$ and $N$ are the
system's own.

For small changes $\dd E$, $\dd V$ and $\dd N$, the entropy changes by one
term for each, its slope times the change (\cref{sec:en-gradient}):
\begin{align*}
   \dd S &= \frac{\dd E}{T} + \frac{P\,\dd V}{T} - \frac{\mu\,\dd N}{T} ,\\
   \dd E &= T\,\dd S - P\,\dd V + \mu\,\dd N
   && \text{multiplying by $T$ and rearranging.}
\end{align*}
$TS$ changes by $T\,\dd S + S\,\dd T$, by the product rule, so
\begin{equation*}
   \dd F = \dd E - T\,\dd S - S\,\dd T = -S\,\dd T - P\,\dd V + \mu\,\dd N .
\end{equation*}
Since $T$ is the slope of $E$ along $S$, subtracting $TS$ exchanged $S$
for that slope as the variable the energy depends on. By
\cref{sec:ha-legendre}, $TS - E$ at the $S$ where the slope is $T$ is the
largest gap between the line $TS$ and the curve $E(S)$, the Legendre
transform of $E$; $F = E - TS$ is that transform with its sign reversed,
and so the least value of $E - TS$ over $S$, which is why the system
settles there. Adding $PV$ or subtracting $\mu N$ exchanges $V$ for $P$ or
$N$ for $\mu$ in the same way (\cref{ex:sm-legendre}).
\Cref{tab:sm-ensembles} collects the four ensembles and what samples each.

\begin{table}[htb]
\centering
\caption{The ensembles of this book: what is fixed, the probability of a
state, the quantity that is greatest or least, and what samples the
ensemble in a simulation.}
\label{tab:sm-ensembles}
\small
\begin{tabular}{@{}lllll@{}}
\toprule
ensemble & fixed & weight & extremum & sampled by\\
\midrule
microcanonical & $N$, $V$, $E$ & equal on the shell & $S$ largest &
   Verlet (Chs~9--11)\\
canonical & $N$, $V$, $T$ & $\mathrm{e}^{-\beta E}$ & $F$ least &
   thermostats (Ch.~13)\\
isothermal-isobaric & $N$, $P$, $T$ & $\mathrm{e}^{-\beta(E + PV)}$ &
   $F + PV$ least & barostats (Ch.~14)\\
grand canonical & $\mu$, $V$, $T$ & $\mathrm{e}^{-\beta(E - \mu N)}$ &
   $F - \mu N$ least & insertion and removal\\
\bottomrule
\end{tabular}
\end{table}

Molecular dynamics keeps the number of atoms, so the grand canonical
ensemble of atoms is sampled by moves that insert or remove an atom at
random and keep or undo each with a probability set by its weight, Monte
Carlo moves, which are not dynamics; this book does not use them. The
grand canonical ensemble returns for electrons.

\subsection*{Electrons in an orbital}

In an electronic-structure calculation the electrons of a solid fill
orbitals, states of definite energy $\varepsilon_{\mathrm o}$ that one electron can
occupy. No two electrons can be in the same state (the exclusion
principle), and an electron has a further two-valued property, its spin,
so an orbital holds either no electron of a given spin or one. In contact
with a reservoir of electrons at the chemical potential $\mu$, the orbital
with one spin is a system of two states: empty, with $N = 0$ and energy
0, and filled, with $N = 1$ and energy $\varepsilon_{\mathrm o}$. Its grand partition
function has two terms, $\Xi = 1 + \mathrm{e}^{-\beta(\varepsilon_{\mathrm o} -
\mu)}$, and its mean occupation, $0\times\mathcal{P}(\text{empty}) +
1\times\mathcal{P}(\text{filled})$, is the probability of the second,
\begin{align*}
   f &= \frac{\mathrm{e}^{-\beta(\varepsilon_{\mathrm o} - \mu)}}{1 +
   \mathrm{e}^{-\beta(\varepsilon_{\mathrm o} - \mu)}}\\
   &= \frac{1}{\mathrm{e}^{(\varepsilon_{\mathrm o} - \mu)/\kB T} + 1}
   && \text{dividing above and below by
   $\mathrm{e}^{-\beta(\varepsilon_{\mathrm o} - \mu)}$,}
\end{align*}
the Fermi-Dirac occupation: 0.5 at $\varepsilon_{\mathrm o} = \mu$, 0.0205 at
\SI{0.1}{\electronvolt} above it and 0.9795 at \SI{0.1}{\electronvolt}
below, at \SI{300}{\kelvin}. The entropy $-\kB\sum\mathcal{P}\ln\mathcal{P}$
of its two states, with probabilities $1 - f$ and $f$, is $-\kB[f\ln f +
(1 - f)\ln(1 - f)]$, whose slope in $f$, $\kB\ln\bigl((1 - f)/f\bigr)$,
vanishes at $f = \frac12$: the entropy is largest, $\kB\ln 2$, when the
orbital is half filled.

\begin{bridge}
Mermin showed that density functional theory carries over to a finite
temperature in the grand canonical ensemble, the grand potential of the
electrons being least at their true density\cite{Mermin1965}. A
calculation of a metal that smears its occupations with this $f$
minimises the free energy $E - TS$ of its electrons, with $S$ the sum of
the orbitals' entropies, and the force on each atom is the slope of that
free energy, so a run on such forces keeps constant the sum of the atoms'
kinetic energy and that free energy\cite{Wentzcovitch1992}, and not the
sum with $E$. A calculation of an electrode held at a fixed potential
fixes $\mu$ of the electrons instead of their number, and makes the grand
potential $E - TS - \mu N$ least: its electrons are in the grand
canonical ensemble, while its atoms may be canonical.
\end{bridge}

\begin{keyidea}
A bath that takes volume weights a state by $\mathrm{e}^{-\beta(E +
PV)}$, and one that takes atoms by $\mathrm{e}^{-\beta(E - \mu N)}$, with
identical atoms counted once, $N!$ orderings as one microstate. In
contact with a bath, the system makes a free energy least: $F = E - TS$ at
fixed $V$, $F + PV$ at fixed $P$, $F - \mu N$ at fixed $\mu$, each the
Legendre transform of the energy with its sign reversed. Fermi-Dirac
occupations are the grand canonical ensemble of one orbital.
\end{keyidea}

\notebook{12}{set a temperature and watch the occupations and their entropy}

\begin{selfcheck}
What is the occupation of an orbital $3\kB T$ below the chemical
potential?
\end{selfcheck}
